\documentclass[11pt,a4paper]{article}
\usepackage[T1]{fontenc}
\usepackage{lmodern,microtype}
\usepackage[margin=27mm,headheight=15pt]{geometry}
\usepackage{amsmath,amssymb,amsthm,mathtools}
\usepackage{xcolor,booktabs,tabularx,longtable,enumitem}
\definecolor{inkblue}{HTML}{244860}
\usepackage[colorlinks=true,linkcolor=inkblue,urlcolor=inkblue]{hyperref}
\usepackage{fancyhdr}
\pagestyle{fancy}\fancyhf{}
\fancyhead[L]{\small Exact symplectic comparison}
\fancyhead[R]{\small Weak taming}
\fancyfoot[C]{\small\thepage}
\setlength{\parskip}{4pt}
\setlength{\parindent}{0pt}
\setlength{\emergencystretch}{2em}
\setlist{itemsep=3pt,topsep=5pt}
\newtheorem{theorem}{Theorem}[section]
\newtheorem{lemma}[theorem]{Lemma}
\newtheorem{proposition}[theorem]{Proposition}
\newtheorem{corollary}[theorem]{Corollary}
\theoremstyle{definition}
\newtheorem{definition}[theorem]{Definition}
\newtheorem{remark}[theorem]{Remark}
\newcommand{\R}{\mathbb R}
\newcommand{\C}{\mathbb C}
\newcommand{\dd}{\mathrm d}
\newcommand{\id}{\mathrm{id}}
\newcommand{\norm}[1]{\left\lVert#1\right\rVert}
\newcommand{\abs}[1]{\left|#1\right|}
\newcommand{\ip}[2]{\left\langle#1,#2\right\rangle}
\newcommand{\flatmap}{\omega^\flat}
\newcommand{\dual}{g^*}
\newcommand{\sref}{\mathcal R}
\allowdisplaybreaks[2]

\hypersetup{pdftitle={Exact symplectic comparison: definitions and target theorems}}
\title{Exact symplectic comparison under weak taming\\[5pt]
\large Definitions, hypotheses, and target theorems}
\author{Exact Liouville comparison project}
\date{24 October 2026}
\begin{document}
\maketitle
\begin{abstract}
We specify an exact comparison theorem for two symplectic primitives on
the same smooth manifold. A radial reference fixes the almost complex
structure and every norm in the hypotheses. The other symplectic form
need only weakly tame that structure and satisfy an integral growth
condition on its inverse. The conclusion includes an actual
diffeomorphism, a smooth isotopy to the identity, and a globally smooth
primitive correction. We also state the finite-type, properly embedded
Stein, and common-reference consequences, with their input data and
quantifier order explicit. A separate construction theorem supplies a
global cylindrical compatible structure and extends arbitrary prescribed
compatible contact data, including the end seam. The companion proof
blueprint supplies the mathematical arguments.
\end{abstract}

\section{Underlying geometry and conventions}
Throughout, $X$ is a finite-dimensional, Hausdorff, second-countable,
smooth real manifold without boundary. Smooth means $C^\infty$.
One fixed real dimension $d\ge0$ is used throughout $X$, with $d$
arbitrary. A nondegenerate two-form forces $d$ to be even when $X$
is nonempty.
The zero-dimensional case is allowed.

A symplectic form is a smooth, closed, pointwise nondegenerate real
two-form. An \emph{exact symplectic structure with specified primitive}
is a smooth real one-form $\lambda$ for which $\omega=\dd\lambda$
is symplectic. The primitive is part of the data: different primitives
of the same two-form are not identified.

Our contraction convention is
\[
 (\iota_v\omega)(w)=\omega(v,w).
\]
Thus the Liouville field of a primitive $\lambda$ is the unique smooth
field $Z_\lambda$ satisfying $\iota_{Z_\lambda}\dd\lambda=\lambda$.
In this document, a \emph{Liouville manifold} has this field complete
for every real time, in both directions. Exact symplectic structures
without this completeness requirement are named as such.

\begin{definition}[Compatibility and the reference metric]
Let $\Omega$ be symplectic. A smooth almost complex structure $J$
is $\Omega$-compatible if
\[
 J^2=-\id,\qquad \Omega(Jv,Jw)=\Omega(v,w),\qquad
 \Omega(v,Jv)>0\quad(v\ne0).
\]
It determines the Riemannian metric
\[
 g(v,w)=\Omega(v,Jw).
\]
\end{definition}
Integrability of $J$ is not required. Neither the definition of
compatibility nor the main theorem assumes completeness of $g$.

\subsection{The tangent, dual, and inverse-form norms}
All norms in the comparison criterion are induced by this particular
$g$. At $x\in X$, write
\[
 |v|_g^2=g_x(v,v),\qquad
 |\xi|_{\dual}=\sup_{|v|_g\le1}|\xi(v)|.
\]
For any symplectic form $\omega$, contraction gives an isomorphism
\[
 \omega_x^\flat:T_xX\longrightarrow T_x^*X,\qquad
 v\longmapsto\omega_x(v,\cdot).
\]
The notation $\|\omega^{-1}\|_g$ means exactly
\begin{equation}\label{eq:inversenorm}
 \|\omega_x^{-1}\|_g
 =\|(\omega_x^\flat)^{-1}\|_{\dual\to g}
 =\sup_{|\xi|_{\dual}\le1}|(\omega_x^\flat)^{-1}\xi|_g .
\end{equation}
In particular, this is not a coordinate-matrix norm chosen independently
of the reference metric. The operator-norm definition also applies in
dimension zero.

\section{The radial reference and admissible growth}
\begin{definition}[Radial reference]\label{def:reference}
A radial reference is a tuple
\[
 \sref=(X,\beta,J,g,R,c,r_*),
\]
where $\beta$ is a smooth real one-form, $\Omega=\dd\beta$ is
symplectic, $J$ is $\Omega$-compatible, and
$g=\Omega(\,\cdot\,,J\,\cdot\,)$.
The function $R:X\to[1,\infty)$ is smooth and proper, and the constants
satisfy $c>0$ and $r_*>1$. On the open set $\{R>r_*\}$ one requires
\begin{equation}\label{eq:radial}
 \beta=-c\,\dd R\circ J,\qquad
 a:=c|\dd R|_{\dual}^2\le R.
\end{equation}
\end{definition}
Here proper means that inverse images of compact sets are compact.
In particular, $\{R\le b\}$ is compact for every finite $b$.
The two identities in \eqref{eq:radial} are imposed only outside
the compact core $\{R\le r_*\}$. No radial identity is imposed on
that core.

\begin{definition}[Admissible growth function]\label{def:growth}
For a lower endpoint $\ell\ge1$, a smooth function
$B:[\ell,\infty)\to(0,\infty)$ is admissible if, for every $r\ge\ell$,
\begin{equation}\label{eq:admissible}
 B(r)\ge r,\qquad
 \lim_{b\to\infty}\int_\ell^b\frac{r}{B(r)^2}\,\dd r=+\infty .
\end{equation}
\end{definition}
Smoothness here means smoothness within the closed half-line, including
its endpoint; values of a representative outside that half-line are
irrelevant. The applications below use $\ell=1$ and $\ell=e$.
The integral condition means divergence to positive infinity of the
ordinary finite-interval integrals; it is not merely a statement
that an unspecified improper integral fails to exist.

\begin{definition}[Weak taming]\label{def:weak}
A symplectic form $\omega$ weakly tames $J$ if
\[
 \omega_x(v,J_xv)\ge0
 \quad\text{for every }x\in X\text{ and }v\in T_xX.
\]
\end{definition}
Nondegeneracy of $\omega$ is a separate assumption. The displayed
inequality is deliberately nonstrict, including for nonzero vectors.

For example, $B(r)=r$ is admissible, as is
$B(r)=r\sqrt{1+\log r}$ on $[1,\infty)$.
The latter example allows a larger inverse-form bound than linear
growth. The theorem below applies to every admissible $B$, not just
to these examples.

\section{The exact comparison theorem}
\begin{definition}[Exact comparison with identity isotopy]\label{def:comparison}
An exact comparison from $\lambda$ to $\beta$ consists of a
diffeomorphism $F:X\to X$, a smooth function $h:X\to\R$, and a smooth
isotopy $(F_t)_{0\le t\le1}$ such that
\[
 F_0=\id,\qquad F_1=F,\qquad
 F^*\lambda=\beta+\dd h,\qquad F^*\dd\lambda=\dd\beta .
\]
Every $F_t$ is a diffeomorphism. Both $(t,x)\mapsto F_t(x)$
and $(t,x)\mapsto F_t^{-1}(x)$ are jointly smooth within
$[0,1]\times X$, including the endpoints. Equivalently, one may
index the diffeomorphisms by all real $t$, with the smoothness and
endpoint requirements imposed only on this closed interval.
\end{definition}
The isotopy is through diffeomorphisms; the statement does not impose
preservation of a fixed symplectic form at every intermediate time.
Neither $h$ nor the isotopy is required to have compact support.

\begin{theorem}[Radial exact comparison]\label{thm:main}
Fix any radial reference $\sref$ as in Definition~\ref{def:reference}.
Let $\lambda$ be any smooth real one-form on $X$ for which
$\omega=\dd\lambda$ is symplectic and weakly tames $J$.
Suppose that there exist an admissible function
$B:[1,\infty)\to(0,\infty)$, a constant $C>0$, and a compact
subset $K\subset X$ such that
\begin{equation}\label{eq:bound}
 \|\omega_x^{-1}\|_g\le C B(R(x))
 \qquad (x\in X\setminus K).
\end{equation}
Then an exact comparison from $\lambda$ to $\beta$ exists:
there are $F$, $h$, and $(F_t)$ as in Definition~\ref{def:comparison},
with
\[
 \boxed{\ F^*\lambda=\beta+\dd h,\qquad F^*\omega=\Omega.\ }
\]
\end{theorem}

\subsection{Quantifier order}
The assertion is uniform over all dimensions, all manifolds as above,
and all radial references:
\[
\begin{gathered}
 \forall\,\sref\quad\forall\,\lambda,\quad
 \left[
 \begin{array}{c}
 \dd\lambda\text{ symplectic and weakly taming},\\
 \exists\,B,C,K\text{ satisfying \eqref{eq:admissible}--\eqref{eq:bound}}
 \end{array}\right]\\
 \Longrightarrow
 \exists\,F,h,(F_t)\text{ with the stated properties}.
\end{gathered}
\]
The witnesses may depend on the reference, the primitive, and the
growth data. The compact set $K$ may differ from the radial core.
Only the estimate \eqref{eq:bound} has this compact-set exception;
positivity and domination of $B$ hold on its entire half-line.

There is no assumption on the growth of $|\lambda|_{\dual}$,
$|\lambda-\beta|_{\dual}$, or their derivatives. There is no
assumed Moser flow or isotopy. Completeness of the metric and of
either primitive's Liouville field is not an input.
The second boxed equality follows from the first by taking exterior
derivatives, but is included to specify the symplectic comparison
explicitly.

\begin{corollary}[Complete Liouville specialization]\label{cor:complete}
If, in addition to the hypotheses of Theorem~\ref{thm:main}, both
specified primitives have complete Liouville fields, the same
conclusion is an exact comparison between those Liouville manifolds.
\end{corollary}
The correction $\dd h$ remains part of the conclusion. The theorem
does not assert $F^*\lambda=\beta$, or conjugacy of the two specified
Liouville fields. Completeness belongs to the supplied primitives;
it is not inferred from a change of primitive by an exact one-form.

\section{A finite-type Liouville reference}
A Liouville domain is given in a boundaryless smooth extension $X_0$
of dimension $d$. Its compact carrier $W\subset X_0$ is locally a
regular sublevel set $\{\rho\le0\}$ at every boundary point. A smooth
$(d-1)$-manifold $Y$ is identified with $\partial W$ by a smooth
embedding $b:Y\to X_0$; if $d=0$, $Y$ is empty. A smooth one-form
$\beta_0$ and smooth vector field $Z_0$ are defined on $X_0$.
On $W$, $\dd\beta_0$ is symplectic and
$\iota_{Z_0}\dd\beta_0=\beta_0$; at the boundary,
$\dd\rho(Z_0)>0$. Put $\alpha=b^*\beta_0$.

Fix a smooth completion realization of this domain. Its data are a
complete Liouville manifold $(X,\beta)$ with flow $\varphi_t$,
a smooth embedding $j:X_0\to X$, and a smooth immersion
$E:Y\times\R\to X$, satisfying
\[
 j^*\beta=\beta_0,\quad \dd j(Z_0)=Z_\beta\circ j
 \quad\hbox{on }W,
\]
\[
 E(y,t)=\varphi_t(j(b(y))),\qquad E^*\beta=e^t\alpha
 \quad(t\in\R).
\]
The restriction of $E$ to $Y\times[0,\infty)$ is a topological
embedding, and
\[
 X=j(W)\cup E(Y\times[0,\infty)),\qquad
 E(Y\times(0,\infty))\cap j(W)=\varnothing .
\]
Thus the completion, its compact-piece inclusion, and its end
parametrization are supplied data. In the end coordinate $r=e^t$,
these data give the customary notation
$X=\widehat W=W\cup_Y(Y\times[1,\infty))$ and $\beta=r\alpha$,
with $W$ identified with $j(W)$.

Let $R_\alpha$ denote the Reeb field, characterized by
$\alpha(R_\alpha)=1$ and $\iota_{R_\alpha}\dd\alpha=0$.
Fix a $\dd\beta$-compatible smooth $J$ which is cylindrical on the end:
\[
 J(r\partial_r)=R_\alpha,\qquad J(\ker\alpha)=\ker\alpha,
\]
and whose restriction to $\ker\alpha$ is independent of $r$.
More explicitly, extend that restriction to $J_Y:TY\to TY$ by
$J_Y R_\alpha=0$. The required end identity is, for $t\ge0$,
\[
 J\,\dd E_{(y,t)}(v,s)
 =\dd E_{(y,t)}\bigl(J_Yv+sR_\alpha,-\alpha(v)\bigr),
 \qquad \alpha(J_Yv)=0.
\]
Set $g=\dd\beta(\,\cdot\,,J\,\cdot\,)$.

\begin{corollary}[Finite-type reference]\label{cor:finite}
For every such reference and $J$, let $\lambda$ be any smooth real
one-form for which $\omega=\dd\lambda$ is symplectic and weakly
tames $J$. If there are an admissible $B:[1,\infty)\to(0,\infty)$,
$C>0$, and $r_0\ge1$ such that
\[
 \|\omega^{-1}_{(y,r)}\|_g\le C B(r)
 \qquad(y\in Y,\ r\ge r_0),
\]
then there exist a diffeomorphism $F$ smoothly isotopic to the identity
and a smooth $h$ satisfying
\[
 F^*\lambda=\beta+\dd h,\qquad F^*\omega=\dd\beta.
\]
\end{corollary}
The bound is equivalently quantified over $t\ge t_0\ge0$ at
$E(y,t)$, with right-hand side $C B(e^t)$ and $r_0=e^{t_0}$.
The new primitive $\lambda$ is not required to be complete or of
finite type. The reference is the specified completion with its
specified end primitive. The radial exhaustion used in applying the
main theorem is constructed from this end; it is not an extra input
to the corollary.

\subsection{Constructing and prescribing the cylindrical structure}
The preceding corollary takes a cylindrical $J$ as input. The following
additional theorem constructs such an input on the specified completion;
it does not add an existence obligation to the general radial theorem.

Write $\xi=\ker\alpha\subset TY$. A prescribed compatible structure on
$\xi$ is a smooth bundle endomorphism $j_\xi$ with
\[
 j_\xi^2=-\id,\qquad
 \dd\alpha(j_\xi v,j_\xi w)=\dd\alpha(v,w),\qquad
 \dd\alpha(v,j_\xi v)>0\quad(0\ne v\in\xi).
\]
These statements use the actual restricted exterior derivative and the
native smooth kernel bundle. They require no integrability.
Put $\pi_\xi(v)=v-\alpha(v)R_\alpha$.

\begin{theorem}[Cylindrical existence and prescribed extension]
\label{thm:cylindrical}
For every Liouville domain and completion realization specified above,
there exists a smooth $\dd\beta$-compatible cylindrical structure $J$
on all of $X$. More precisely, for every prescribed smooth compatible
$j_\xi$, there exists such a $J$ satisfying, for every $y\in Y$, $t\ge0$,
$v\in T_yY$, and $s\in\R$,
\begin{equation}\label{eq:prescribed-end}
 J_{E(y,t)}\bigl(\dd E_{(y,t)}(v,s)\bigr)
 =\dd E_{(y,t)}\bigl(j_\xi\pi_\xi(v)+sR_\alpha,
                         -\alpha(v)\bigr).
\end{equation}
In particular, equality holds at $t=0$, the seam $r=1$, as well as on
the entire positive end. The output obeys $J^2=-\id$ and determines
the smooth positive-definite metric
\[
 g(v,w)=\dd\beta(v,Jw).
\]
The existence of at least one smooth compatible $j_\xi$ is also proved;
it is not an assumption of the first existence assertion.
\end{theorem}

The two assertions have different quantifiers: existence chooses some
contact data, whereas prescribed extension must retain every allowed
input $j_\xi$ exactly. No positive-rank or nonempty-boundary hypothesis
is added. The permitted zero-dimensional and empty cases remain included.

\begin{corollary}[A constructed finite-type reference]
For every specified domain and completion, one can choose $J$ and its
metric $g$ before choosing the comparison primitive $\lambda$, so that
every $\lambda$ satisfying the weak taming and eventual growth
hypotheses of Corollary~\ref{cor:finite} has the stated exact comparison.
The same assertion applies to a $J$ extending prescribed $j_\xi$.
\end{corollary}
The target hypotheses remain relative to this chosen $J$ and its
explicit metric. Constructing $J$ does not assert that every symplectic
form weakly tames it. The separate native finite-type presentation
adapter constructs the completion from its own source data; the
integrated comparison uses that same quotient, primitive, and end map
throughout, rather than identifying independently chosen completions.

\clearpage
\section{A properly embedded Stein reference}
Let $X$ be a smooth real manifold as above, and let $J:TX\to TX$
be a smooth bundle endomorphism with $J^2=-\id$. Suppose we are given
a smooth proper topological embedding
\[
 \iota:X\hookrightarrow\C^N.
\]
Require its actual differential to be injective at every point and
to satisfy the Cauchy--Riemann identity
\[
 \dd\iota_x(J_xv)=I\,\dd\iota_x(v),\qquad Iu=iu.
\]
These are the precise input data used here. They include the usual
case of a proper holomorphic embedding of a complex manifold; an
independent complex atlas or integrability certificate is not an
additional hypothesis. The embedding is part of the data.
Define, using its actual
differential and the standard Euclidean norm,
\begin{equation}\label{eq:stein}
 R=e+\|\iota\|^2,\qquad
 \beta=-\frac14\,\dd R\circ J,\qquad \Omega=\dd\beta .
\end{equation}

\begin{proposition}[The reference determined by the embedding]\label{prop:stein}
These formulas determine a smooth exact symplectic structure
compatible with the given $J$. Its metric is exactly the induced
Euclidean metric:
\[
 g_x(v,w)=\Omega_x(v,Jw)
   =\ip{\dd\iota_x(v)}{\dd\iota_x(w)}_{\R}.
\]
The function $R$ is smooth and proper, $R\ge e$, and
\[
 |\dd R|_{\dual}^2\le4(R-e),\qquad
 \frac14|\dd R|_{\dual}^2\le R-e<R.
\]
\end{proposition}
The real inner product here is the real part of the Hermitian
Euclidean product. Proposition~\ref{prop:stein} verifies the
reference formulas from the embedding; a symplectic form or induced
smooth metric with these identities is not supplied as an additional
hypothesis.

\begin{corollary}[Stein reference]\label{cor:stein}
For every embedding and almost complex structure specified above,
form the determined
reference \eqref{eq:stein}. For every smooth real one-form $\lambda$
such that $\omega=\dd\lambda$ is symplectic and weakly tames $J$,
assume that there exist $C>0$, a compact $K\subset X$, and an
admissible $B:[e,\infty)\to(0,\infty)$ such that
\[
 \|\omega_x^{-1}\|_g\le C B(R(x))\qquad(x\notin K).
\]
Then there are a diffeomorphism $F$ smoothly isotopic to the identity
and a smooth $h$ such that
\[
 F^*\lambda=\beta+\dd h,\qquad F^*\omega=\Omega.
\]
\end{corollary}
No finite-type hypothesis on the Stein manifold is imposed.
The growth condition is measured in the metric of the chosen
embedding. The reference is determined before the universally
quantified comparison primitive $\lambda$ is chosen.

\clearpage
\section{Comparison through one common reference}
\begin{corollary}[Common radial reference]\label{cor:common}
Fix a single radial reference $\sref$.
Let $\lambda_1,\lambda_2$ each satisfy the assumptions of
Theorem~\ref{thm:main} relative to that reference.
Their growth functions $B_i$, constants $C_i$, and compact
exceptional sets $K_i$ may be different.
Then there are a diffeomorphism $\Phi:X\to X$, smoothly isotopic
to the identity, and a smooth function $q:X\to\R$ such that
\[
 \Phi^*\lambda_1=\lambda_2+\dd q,\qquad
 \Phi^*\dd\lambda_1=\dd\lambda_2.
\]
\end{corollary}
If $F_i^*\lambda_i=\beta+\dd h_i$ are the two comparisons, the
required witnesses are explicitly
\[
 \Phi=F_1\circ F_2^{-1},\qquad
 q=(h_1-h_2)\circ F_2^{-1}.
\]
Thus the common reference fixes the geometric standard against
which the two independent growth conditions are tested.

\section{Scope of the conclusion}
The statements concern exact symplectic identifications of the
whole open manifolds, together with smooth identity isotopies.
They do not impose compact support, preservation of the primitives
without an exact correction, or a controlled finite-type
deformation of all intermediate primitives.
The term Liouville is used with all-real-time completeness, whereas
the general comparison theorem applies without that assumption.

The companion \emph{Proof blueprint} constructs four Moser evolutions,
proves existence in both parameter directions, derives the exact
primitive correction, constructs the cylindrical extension, and verifies
the two geometric reference applications above. The comparison and the isotopy are conclusions
of that argument, not fields assumed in the input data.
\end{document}
